\documentclass[12pt]{article}
\usepackage[T1]{fontenc}
\usepackage[utf8]{inputenc}
\usepackage{lmodern}
\usepackage{textcomp}
\usepackage{enumitem}
\usepackage{geometry}
\geometry{margin=1in}
\begin{document}
\begin{center}
Answer Key - Physics - Undergraduate Level Assessment 6 \
\small (Answers are ordered exactly as in the question paper.)
\end{center}

\section*{Multiple Choice}
\begin{enumerate}[label=\arabic*.]
\item \textbf{Answer:} A
\\ \textit{Options:} A) 4; B) 5; C) 6; D) 20
\\ \textit{Note:} The angular magnification of a refracting telescope is determined by the ratio of the focal lengths of the objective lens to the eyepiece lens, and given that the focal length of the eyepiece is 20 cm, the correct angular magnification is 4 when the focal length of the objective lens is 80 cm, resulting in a magnification of 80/20 = 4.
\item \textbf{Answer:} D
\\ \textit{Options:} A) Decreases by a factor of 81.; B) Decreases by a factor of 9.; C) Increases by a factor of 9.; D) Increases by a factor of 81.
\\ \textit{Note:} The energy radiated per second per unit area from a blackbody is described by the Stefan-Boltzmann law, which states that it is proportional to the fourth power of the absolute temperature, so if the temperature is increased by a factor of 3, the energy increases by a factor of \(3^4 = 81\).
\item \textbf{Answer:} D
\\ \textit{Options:} A) the weak nuclear force; B) the strong nuclear force; C) vacuum polarization; D) interactions with the ionic lattice
\\ \textit{Note:} The correct answer is due to the fact that in BCS theory, Cooper pairs form as electrons interact with the lattice vibrations (phonons) of the ionic lattice, leading to an effective attractive force that allows for superconductivity.
\item \textbf{Answer:} D
\\ \textit{Options:} A) 0.1 GeV/$c^2$; B) 0.2 GeV/$c^2$; C) 0.5 GeV/$c^2$; D) 1.0 GeV/$c^2$
\\ \textit{Note:} The rest mass of a particle can be calculated using the equation \(E^2 = (pc)^2 + (m_0 c^2)^2\), and substituting the given total energy of 5.0 GeV and momentum of 4.9 GeV/c yields a rest mass of approximately 1.0 GeV/c².
\item \textbf{Answer:} A
\\ \textit{Options:} A) Electron; B) Meson; C) Photon; D) Boson
\\ \textit{Note:} The negative muon (mu^-) has properties most similar to the electron because both are leptons, share a similar charge of -1, and have comparable behavior in electromagnetic interactions, despite the muon's greater mass.
\end{enumerate}

\section*{True / False}
\begin{enumerate}[label=\arabic*.]
\setcounter{enumi}{5}
\item \textbf{Answer:} False
\\ \textit{Options:} A) True; B) False
\\ \textit{Note:} The statement is false because the n = 1 level can hold a maximum of 2 electrons, while the n = 2 level can hold a maximum of 8 electrons, resulting in a total of 10 electrons for both levels when filled.
\item \textbf{Answer:} True
\\ \textit{Options:} A) True; B) False
\\ \textit{Note:} For an electron in the l = 2 state, the magnetic quantum number m\_l can take on integer values ranging from -l to +l, which in this case gives the five allowed values of -2, -1, 0, +1, and +2.
\item \textbf{Answer:} True
\\ \textit{Options:} A) True; B) False
\\ \textit{Note:} The statement is true because the relativistic time dilation effect allows the particle, which has a proper lifetime that would otherwise be too short to travel 450 m at 0.60 c, to appear to travel that distance in the lab frame due to its increased time experienced by an outside observer.
\item \textbf{Answer:} True
\\ \textit{Options:} A) True; B) False
\\ \textit{Note:} The angular magnification of a refracting telescope is given by the ratio of the focal lengths of the objective lens to the eyepiece lens, and with a 20 cm focal length eyepiece and a suitable objective lens, the magnification can indeed be 4.
\item \textbf{Answer:} True
\\ \textit{Options:} A) True; B) False
\\ \textit{Note:} The resolving power of a grating spectrometer is defined as the ratio of the wavelength to the minimum difference in wavelength that can be resolved, and since a resolving power of 250 indicates that it can distinguish between wavelengths that differ by 2 nm (from 500 nm to 502 nm), the statement is true.
\end{enumerate}

\section*{Long Form Response}
\begin{enumerate}[label=\arabic*.]
\setcounter{enumi}{10}
\item \textbf{Answer:} Characteristic X-rays are generated when high-energy electrons collide with a metal target, resulting in the ejection of inner shell electrons from the metal atoms. This creates vacancies in the inner electron shells, specifically in the K or L shells, depending on the energy of the incident electrons. As outer shell electrons transition to fill these vacancies, they release energy in the form of X-rays, which have specific energies characteristic of the metal used in the target. This process not only highlights the interaction between high-energy electrons and metal atoms but also underscores the significance of inner shell electron vacancies in the production of characteristic radiation. Thus, the emission of X-rays is fundamentally linked to the filling of these inner shell vacancies by outer shell electrons.
\\ \textit{Note:} The answer correctly describes the generation of characteristic X-rays by detailing how high-energy electron collisions eject inner shell electrons, creating vacancies that are subsequently filled by outer shell electrons, resulting in the emission of X-rays with energies specific to the metal target.
\item \textbf{Answer:} Bosons and fermions are two fundamental categories of particles in quantum physics, distinguished primarily by their wave function symmetries. Fermions, which include particles like electrons, protons, and neutrons, possess antisymmetric wave functions, meaning that exchanging two identical fermions results in a change of sign of the wave function. This property leads to the Pauli exclusion principle, which states that no two fermions can occupy the same quantum state simultaneously, resulting in phenomena such as electron shell structure in atoms and the stability of matter. In contrast, bosons, such as photons and gluons, have symmetric wave functions, allowing multiple bosons to occupy the same state, which underlies phenomena like Bose-Einstein condensation. Thus, the distinctions in wave function symmetries fundamentally influence the collective behaviors and interactions of these particles in various physical systems.
\\ \textit{Note:} correct because it accurately describes the fundamental distinction between bosons and fermions based on their wave function symmetries, highlighting how fermions' antisymmetric wave functions give rise to the Pauli exclusion principle, which dictates their behavior in quantum systems.
\end{enumerate}

\vfill
\noindent\textit{End of Answer Key}
\end{document}